
\clearpage
\VerdierPhysicalPageAnchor{0128}%
\begingroup
\small
\setlength{\emergencystretch}{1.5em}%
\setlength{\abovedisplayskip}{3pt plus 2pt minus 2pt}%
\setlength{\belowdisplayskip}{3pt plus 2pt minus 2pt}%

\VerdierSemanticAnchor{ch2-2}
\section{三角圏における局所化。}

\VerdierSemanticAnchor{ch2-2-1}
\subsection{射の乗法的系 [16]。}

\VerdierSemanticAnchor{ch2-2-1-1}
\noindent\textbf{2.1.1.}\quad
\VerdierSemanticAnchor{ch2-2-1-1-multiplicative-system}
$C$ を圏（加法的であるとは限らない）とする。$C$ の射の集合 $S$ が次の
性質をもつとき、これを\emph{乗法的系}と呼ぶ。

\VerdierSemanticAnchor{ch2-2-1-1-sm1}
\smallskip
\noindent\textbf{SM1)} $S$ に属する合成可能な二つの射の合成は $S$ に
属する。$C$ の各対象 $X$ に対し、$X$ の恒等射は $S$ に属する。

\VerdierSemanticAnchor{ch2-2-1-1-sm2}
\smallskip
\noindent\textbf{SM2)} 任意の図式

\VerdierSemanticAnchor{ch2-2-1-1-sm2-right-input}
\begin{center}
  \begin{tikzcd}[column sep=4.1em,row sep=2.5em]
    & {} \arrow[d,"s\in S"] \\
    {} \arrow[r,"f"'] & {}
  \end{tikzcd}
\end{center}
\noindent は次の可換図式に補完できる：

\VerdierSemanticAnchor{ch2-2-1-1-sm2-right-completion}
\begin{center}
  \begin{tikzcd}[column sep=4.1em,row sep=2.5em]
    {} \arrow[r,dashed,"g"] \arrow[d,dashed,"t\in S"']
      & {} \arrow[d,"s\in S"] \\
    {} \arrow[r,"f"'] & {}
  \end{tikzcd}
  \quad .
\end{center}
\noindent 任意の図式

\VerdierSemanticAnchor{ch2-2-1-1-sm2-left-input}
\begin{center}
  \begin{tikzcd}[column sep=4.1em,row sep=2.5em]
    {} \arrow[r,"g"] \arrow[d,"t\in S"'] & {} \\
    {} & {}
  \end{tikzcd}
\end{center}
\noindent は次の可換図式に補完できる：

\VerdierSemanticAnchor{ch2-2-1-1-sm2-left-completion}
\begin{center}
  \begin{tikzcd}[column sep=4.1em,row sep=2.5em]
    {} \arrow[r,"g"] \arrow[d,"t\in S"']
      & {} \arrow[d,dashed,"s\in S"] \\
    {} \arrow[r,dashed,"f"'] & {}
  \end{tikzcd}
  \quad .
\end{center}
\VerdierSemanticAnchor{ch2-2-1-1-sm3-boundary}
\noindent\textbf{SM3)} 任意の二射 $f,g\colon X\rightrightarrows Y$ に
対し、次の性質は同値である。
\endgroup
